\documentclass[aspectratio=169]{beamer}

\usepackage{listings}
\usepackage{xcolor}

\definecolor{codegray}{rgb}{0.5,0.5,0.5}
\definecolor{codepurple}{rgb}{0.58,0,0.82}
\definecolor{backcolour}{rgb}{0.95,0.95,0.92}

\lstdefinestyle{customc}{
  backgroundcolor=\color{backcolour},
  commentstyle=\color{codegray}\itshape,
  keywordstyle=\color{blue}\bfseries,
  numberstyle=\tiny\color{gray},
  stringstyle=\color{codepurple},
  basicstyle=\ttfamily\footnotesize,
  breakatwhitespace=false,
  breaklines=true,
  captionpos=b,
  keepspaces=true,
  numbers=left,
  numbersep=5pt,
  showspaces=false,
  showstringspaces=false,
  showtabs=false,
  tabsize=2,
  frame=single,
  language=C++
}

\lstset{style=customc}

\usepackage[spanish]{babel}
\usepackage[utf8]{inputenc}
\usepackage[T1]{fontenc}
\usepackage{lmodern}
\usepackage{amsmath, amsfonts}
\usepackage{hyperref}

\usetheme{Madrid}

\title{Lógica de Hoare}
\author{Sebastian Mestre}
\date{}

\begin{document}

\begin{frame}
  \titlepage
\end{frame}

\section{Trípletas de Hoare}

\begin{frame}{¿Qué es una trípleta de Hoare?}
  \begin{itemize}
    \item La \textbf{lógica de Hoare} es una forma de razonar formalmente sobre la corrección de programas.
    \item Una \textbf{trípeta de Hoare} se escribe:
      \[
        \{P\}\ C\ \{Q\}
      \]
      donde:
      \begin{itemize}
        \item \(P\): \textbf{precondición}, afirmación lógica sobre el estado antes de ejecutar el código.
        \item \(C\): \textbf{comando} o fragmento de código.
        \item \(Q\): \textbf{postcondición}, afirmación lógica sobre el estado después de ejecutar \(C\).
      \end{itemize}
    \item Lectura: \textbf{“Si \(P\) es verdadera antes de ejecutar \(C\) y se ejecuta \(C\), entonces \(Q\) será verdadera al finalizar”}.
  \end{itemize}
\end{frame}

\begin{frame}[fragile]{Ejemplo simple}
  \begin{itemize}
    \item Consideremos el fragmento:
  \end{itemize}
  \begin{block}{Programa}
  \begin{lstlisting}[style=customc]
x = x * 6;
  \end{lstlisting}
  \end{block}
  \begin{itemize}
    \item El mismo programa puede aparecer en \textbf{muchas} trípletas válidas, según la propiedad que nos interese:
  \end{itemize}
  \begin{block}{Algunos ejemplos}
  \[
    \{\, x \% 2 == 1 \,\}\ x = x * 6;\ \{\, x \% 2 == 1 \,\}
  \]
  \[
    \{\, x \% 2 == 0 \,\}\ x = x * 6;\ \{\, x \% 2 == 0 \,\}
  \]
  \[
    \{\, x \% 2 == 0 \,\}\ x = x * 6;\ \{\, x \% 4 == 0 \,\}
  \]
  \[
    \{\, x > 16 \,\}\ x = x * 6;\ \{\, x > 96 \,\}
  \]
  \end{block}
\end{frame}

\begin{frame}{Precondición más débil}
  \begin{itemize}
    \item Fijada una \textbf{postcondición} \(Q\) y un programa \(C\), hay muchas precondiciones \(P\) tales que:
      \[
        \{P\}\ C\ \{Q\}
      \]
    \item La \textbf{precondición más débil} es aquella \(P_{\text{mín}}\) que es implicada por cualquier otra precondición que garantice \(Q\).
    \item Intuitivamente: es la condición “mínima” que asegura que, tras ejecutar el programa, la postcondición se cumple.
  \end{itemize}
  \begin{block}{Ejemplo}
    Para el programa \(x = x * 6\) y la postcondición
    \[
      x \% 12 == 0,
    \]
    una precondición más débil adecuada es:
    \[
      x \% 2 == 0.
    \]
    Entonces:
    \[
      \{\, x \% 2 == 0 \,\}\ x = x * 6;\ \{\, x \% 12 == 0 \,\}.
    \]
  \end{block}
\end{frame}

\begin{frame}{Cómo se usa en la práctica}
  \begin{itemize}
    \item En teoría existen reglas para calcular (semi-automáticamente) la precondición más débil de una postcondición dada.
    \item En la práctica:
      \begin{itemize}
        \item Se arranca desde la \textbf{postcondición deseada} al final del programa.
        \item Se van propagando hacia atrás las \textbf{condiciones necesarias} para que esa postcondición se cumpla.
        \item Si en algún punto podemos garantizar una precondición \textbf{más fuerte} que la mínima requerida, no hay problema.
      \end{itemize}
    \item El concepto dual de \textbf{postcondición más fuerte} existe, pero es mucho más difícil de calcular y casi no se usa.
  \end{itemize}
\end{frame}

\section{Reglas de composición}

\begin{frame}{Composición secuencial}
  \begin{itemize}
    \item Una regla básica de la lógica de Hoare es la \textbf{composición secuencial}.
  \end{itemize}
  \begin{block}{Regla}
    Si
    \[
      \{P\}\ C_1\ \{Q\}
      \quad\text{y}\quad
      \{Q\}\ C_2\ \{R\},
    \]
    entonces
    \[
      \{P\}\ C_1; C_2\ \{R\}.
    \]
  \end{block}
  \begin{itemize}
    \item Interpretación informal:
      \begin{itemize}
        \item \(C_1\) transforma estados con \(P\) en estados con \(Q\).
        \item \(C_2\) transforma estados con \(Q\) en estados con \(R\).
        \item Por lo tanto, ejecutar primero \(C_1\) y luego \(C_2\) transforma estados con \(P\) en estados con \(R\).
      \end{itemize}
  \end{itemize}
\end{frame}

\section{Bucles e invariantes}

\begin{frame}{Bucle \texttt{while} e invariantes}
  \begin{itemize}
    \item Para razonar sobre un bucle \texttt{while} usamos un \textbf{invariante de bucle} \(I\).
    \item Un invariante es una propiedad lógica que:
      \begin{itemize}
        \item Es verdadera antes de entrar al bucle.
        \item Se mantiene tras cada iteración.
        \item Sigue siendo verdadera cuando el bucle termina.
      \end{itemize}
    \item Además, cuando el bucle termina, la \textbf{condición del while es falsa}.
  \end{itemize}
\end{frame}

\begin{frame}[fragile]{Esquema de razonamiento para \texttt{while}}
  \begin{block}{Esquema}
  \begin{lstlisting}[style=customc]
// {I}
while (condicion) {
  // {I && condicion}
  // ...codigo que debe preservar I...
  // {I}
}
// {I && !condicion}
  \end{lstlisting}
  \end{block}
  \begin{itemize}
    \item Para demostrar la corrección:
      \begin{enumerate}
        \item Encontrar un invariante \(I\) que sea cierto antes de entrar al bucle.
        \item Probar que si \(I\) es cierto al inicio de una iteración y la condición es verdadera, entonces tras el cuerpo del bucle \(I\) sigue siendo cierto.
        \item Al terminar el bucle, usar \(I \land \neg\text{condición}\) para implicar la postcondición deseada.
      \end{enumerate}
  \end{itemize}
\end{frame}

\section{En la práctica competitiva}

\begin{frame}{Lógica de Hoare en la vida real}
  \begin{itemize}
    \item En programación competitiva \textbf{no} escribimos demostraciones formales en notación de Hoare.
    \item Pero los conceptos fundamentales sí se usan todo el tiempo:
      \begin{itemize}
        \item \textbf{Invariantes} de estructuras de datos y bucles.
        \item \textbf{Precondiciones} de funciones y rutinas auxiliares.
        \item \textbf{Postcondiciones} que queremos garantizar al terminar un algoritmo.
      \end{itemize}
    \item Ejemplos típicos:
      \begin{itemize}
        \item “Al inicio de cada iteración, los primeros \(k\) elementos están ordenados”.
        \item “Al entrar al bucle, \texttt{res} contiene el resultado parcial correcto hasta ahora”.
      \end{itemize}
  \end{itemize}
\end{frame}

\begin{frame}{Más allá de los bucles}
  \begin{itemize}
    \item El razonamiento con invariantes se usa también en:
      \begin{itemize}
        \item Recursiones: qué “promesa” cumple cada llamada recursiva.
        \item Ecuaciones y desigualdades en problemas matemáticos.
        \item Recorridos no triviales sobre arreglos y grafos.
      \end{itemize}
    \item Muchas veces “sentimos” que un algoritmo es correcto porque identificamos claramente:
      \begin{itemize}
        \item Qué se mantiene invariante.
        \item Qué precondiciones necesita cada paso.
        \item Qué postcondición queremos al final.
      \end{itemize}
  \end{itemize}
\end{frame}

\section*{Cierre}

\begin{frame}{Ideas clave}
  \begin{itemize}
    \item La lógica de Hoare da un \textbf{lenguaje formal} para hablar de precondiciones, postcondiciones e invariantes.
    \item Aunque no se use la notación completa en concursos, la forma de pensar sí está profundamente influenciada por estas ideas.
    \item Adoptar este estilo de razonamiento es clave para escribir algoritmos correctos y poder justificar por qué funcionan.
  \end{itemize}
\end{frame}

\begin{frame}{Preguntas}
  \centering
  \Huge
  \textbf{Q\&A}
\end{frame}

\end{document}


